\documentclass[11pt]{article}

\usepackage[T1]{fontenc}
\usepackage{amsmath,amssymb,amsthm}
\usepackage{hyperref}

\title{Goodness-of-Fit and Nested Tests for\\ Non-Iterative CFA Estimators}
\author{}
\date{\today}

\newcommand{\E}{\mathbb{E}}
\newcommand{\tr}{\operatorname{tr}}
\newcommand{\Var}{\operatorname{Var}}
\newcommand{\Cov}{\operatorname{Cov}}
\newcommand{\vech}{\operatorname{vech}}
\newcommand{\vecop}{\operatorname{vec}}
\newcommand{\diag}{\operatorname{diag}}
\newcommand{\rank}{\operatorname{rank}}
\newcommand{\IF}{\operatorname{IF}}
\newcommand{\argmin}{\operatorname*{arg\,min}}
\newcommand{\plim}{\operatorname*{plim}}
\newcommand{\T}{^{\mathsf T}}
\newcommand{\Ninv}{^{-1}}

\theoremstyle{plain}
\newtheorem{theorem}{Theorem}
\newtheorem{proposition}{Proposition}
\newtheorem{corollary}{Corollary}
\theoremstyle{definition}
\newtheorem{assumption}{Assumption}
\theoremstyle{remark}
\newtheorem*{remark}{Remark}

\begin{document}
\maketitle

\section*{Purpose}

Non-iterative CFA estimators (Guttman's multiple-group method
\cite{guttman1952} and the wider family evaluated by Dhaene and Rosseel
\cite{dhaene2024}) are consistent, robust, convergence-free, and cheap. Their
standard errors follow by the delta method (companion note
\texttt{guttman\_cfa\_asymptotics}). The one remaining inferential gap is a
\emph{goodness-of-fit} test and a \emph{nested} (difference) test. Both are
supplied by an old and standard idea: evaluate a fixed-weight discrepancy at the
estimates and correct its reference distribution for the fact that the estimates
do not minimize it \cite{browne1984,shapiro1986,satorra1989}. Nothing here is
new asymptotics; the only estimator-specific ingredient is the map's Jacobian,
which the SE work already produces.

The note develops the theory estimator-agnostically first, then specialises to
two discrepancies: unweighted least squares (ULS, the natural and cheapest
choice) and normal-theory maximum likelihood (NTML). We deliberately exclude
GLS: its weight is the \emph{sample} object $\widehat S\Ninv$, and we want the
discrepancy weight to be either fixed (ULS) or a smooth function of the fitted
model (NTML), never a function of the data. The point estimator needs no fourth
moments; every test does, and we take the fourth-moment matrix $\Gamma$
empirical by default. Normality is a special case that simplifies the algebra
but never collapses the reference distribution to a plain $\chi^2$, because these
estimators are not efficient. That is deliberate: efficiency is not the target.
Robustness is, and the downstream target is categorical data, where a
non-iterative estimator on polychoric moments with a distribution-free test is
exactly the underserved cell.

\section{Setup and estimator family}\label{sec:setup}

Let $y_i\in\mathbb R^p$ have covariance $\Sigma_0$ and finite fourth moments.
Write $p^\star=p(p+1)/2$, $\hat\sigma=\vech(\widehat\Sigma)$,
$\sigma_0=\vech(\Sigma_0)$, and the covariance influence rows
$z_i=\vech\{(y_i-\mu_0)(y_i-\mu_0)\T-\Sigma_0\}$, so that
\begin{equation}\label{eq:clt}
  \sqrt N(\hat\sigma-\sigma_0)=N^{-1/2}\sum_i z_i+o_p(1)
  \Rightarrow N(0,\Gamma),\qquad \Gamma=\Var(z_i).
\end{equation}
Under normality $\Gamma$ reduces to
$\Gamma_{\mathrm N}=2\,D^{+}(\Sigma_0\otimes\Sigma_0)D^{+\T}$, with $D$ the
$p^2\times p^\star$ duplication matrix and $D^{+}=(D\T D)\Ninv D\T$; otherwise
$\Gamma$ is the actual fourth-moment covariance and is estimated by its sample
analogue. A CFA model supplies $\sigma(\theta)=\vech\{\Sigma(\theta)\}$ with
$q$ free parameters, model Jacobian $\Delta=\partial\sigma/\partial\theta\T\in
\mathbb R^{p^\star\times q}$ of full column rank $q$, and residual degrees of
freedom $df=p^\star-q$.

A \emph{non-iterative estimator} is a fixed closed-form map
$\hat\theta=\tau(\hat\sigma)$.

\begin{assumption}\label{ass:reg}
$\tau$ is Fisher-consistent and differentiable at $\sigma_0$: for every $\theta$
in the model, $\tau\{\sigma(\theta)\}=\theta$, and $J=D\tau(\sigma_0)\in
\mathbb R^{q\times p^\star}$ exists. Fisher consistency gives the identity
\begin{equation}\label{eq:JDelta}
  J\Delta=I_q .
\end{equation}
\end{assumption}

Equation \eqref{eq:JDelta} is the one structural fact the test theory uses. It
holds for the estimators below whenever their communality or instrument rule is
exact on the model manifold; Dhaene and Rosseel's near-zero bias confirms it
empirically.

\paragraph{Members.} All are measurement-model (CFA) estimators; full SEM is out
of scope, since the structural part needs a separate (SAM-style two-step or
joint) treatment.
\begin{itemize}
\item \emph{Multiple-group / Guttman} \cite{guttman1952}. With a $0/1$
  indicator-by-factor pattern $X$, form $C(\sigma)$ by placing communalities on
  the diagonal, then
  $\Lambda_G=CX(X\T CX)\Ninv D_B^{1/2}M\Ninv$,
  $\Phi_G=M(D_B^{-1/2}X\T CXD_B^{-1/2})M$,
  $\psi_G=\diag\{S-\Lambda_G\Phi_G\Lambda_G\T\}$, with $D_B=\diag(X\T CX)$ and
  $M$ the marker rescaling. The full map, its differential, and $J_G$ are in the
  companion note; it returns a complete $(\Lambda,\Phi,\psi)$, not just an
  extraction.
\item \emph{FABIN2/FABIN3} \cite{hagglund1982}: factor analysis by instrumental
  variables; per-indicator closed form using sibling indicators as instruments.
\item \emph{Bentler (1982) noniterative ULS} and a \emph{James--Stein} shrinkage
  variant, as in \cite{dhaene2024}.
\item \emph{MIIV-2SLS} \cite{bollen1996}: model-implied instrumental variables,
  two-stage least squares.
\end{itemize}
Each supplies its own $(\tau,J)$; the theory below is written for a generic
pair. Dhaene and Rosseel's headline is that within this family the choice barely
matters, so the natural default is the cheapest, Guttman.

\section{Residual-based fit test for a general estimator}\label{sec:gof}

Fix a symmetric weight $V\in\mathbb R^{p^\star\times p^\star}$, positive
semidefinite on the relevant subspace, and define the discrepancy statistic
\begin{equation}\label{eq:TV}
  T_V=N\,\hat r\T V\hat r,\qquad \hat r=\hat\sigma-\sigma(\hat\theta).
\end{equation}
Since $\hat\theta$ need not minimise $\hat r\T V\hat r$, the classical
projector does not apply; the estimator's own Jacobian takes its place.

\begin{proposition}[Residual law]\label{prop:gof}
Under Assumption~\ref{ass:reg}, \eqref{eq:clt}, and the model
$\Sigma_0=\Sigma(\theta_0)$,
\begin{equation}\label{eq:res-lin}
  \sqrt N\,\hat r=M\,\sqrt N(\hat\sigma-\sigma_0)+o_p(1),
  \qquad M=I-\Delta J,
\end{equation}
where $M$ is an oblique projector ($M\Delta=0$, $\rank M=df$). Consequently
\begin{equation}\label{eq:weighted-chi2}
  T_V\Rightarrow\sum_{j=1}^{df}\lambda_j\,\chi^2_{1,j},
\end{equation}
where the $\lambda_j$ are the $df$ nonzero eigenvalues of $U_V\Gamma$ with
$U_V=M\T V M$, equivalently of $VM\Gamma M\T$. Its mean is
$\E\,T_V\to\tr(U_V\Gamma)$.
\end{proposition}

\begin{proof}
$\hat r=\hat\sigma-\sigma_0-\Delta(\hat\theta-\theta_0)+o_p(N^{-1/2})$ and
$\hat\theta-\theta_0=J(\hat\sigma-\sigma_0)+o_p(N^{-1/2})$ give
\eqref{eq:res-lin}. $M\Delta=\Delta-\Delta J\Delta=\Delta-\Delta=0$ by
\eqref{eq:JDelta}, so $\Delta J$ is idempotent of rank $q$ and $M$ has rank
$df$. Then $\sqrt N\hat r\Rightarrow M\zeta$, $\zeta\sim N(0,\Gamma)$, and
$T_V\Rightarrow\zeta\T M\T V M\zeta$, a quadratic form in a normal vector whose
law is \eqref{eq:weighted-chi2} with the stated eigenvalues.
\end{proof}

To turn \eqref{eq:weighted-chi2} into a usable reference, reduce the spectrum by
any of the standard devices, all already implemented for scaled statistics in
\texttt{robust}:
\begin{itemize}
\item mean scaling (Satorra--Bentler \cite{satorrabentler1994}):
  $\bar T_V=T_V/\hat c$, $\hat c=\tr(\widehat U_V\widehat\Gamma)/df$, referred
  to $\chi^2_{df}$;
\item mean-and-variance (Satterthwaite): scale and set
  $\hat d=(\sum_j\lambda_j)^2/\sum_j\lambda_j^2$;
\item exact eigenvalue reference (the repo's pEBA/FMG reducer applied to
  $\widehat U_V\widehat\Gamma$).
\end{itemize}

The efficient special case explains what is being given up.

\begin{proposition}[Efficient collapse]\label{prop:collapse}
If in addition $\hat\theta$ is the minimiser of $\hat r\T V\hat r$, so
$J=(\Delta\T V\Delta)\Ninv\Delta\T V$ and $M=M_V:=I-\Delta(\Delta\T
V\Delta)\Ninv\Delta\T V$, and if $V$ is a generalised inverse of $\Gamma$ on the
symmetric subspace ($V\Gamma V=V$), then every $\lambda_j=1$ and
$T_V\Rightarrow\chi^2_{df}$.
\end{proposition}

Both conditions fail for a non-iterative $\tau$: its $J$ is not the
$V$-projector, so $M\neq M_V$, and neither ULS nor NTML uses $V=\Gamma\Ninv$
with the matching estimator. Hence \eqref{eq:weighted-chi2} does not collapse,
and scaling is mandatory rather than cosmetic.

\section{The two discrepancies}\label{sec:disc}

\paragraph{ULS.} With $F_{\mathrm U}(\theta)=\tfrac12\tr[\{\widehat\Sigma-
\Sigma(\theta)\}^2]=\tfrac12\hat r\T V_{\mathrm U}\hat r$ and
$V_{\mathrm U}=D\T D$ (the vech metric: unit weight on variances, weight $2$ on
covariances), $T_{\mathrm U}=2NF_{\mathrm U}(\hat\theta)=N\hat r\T V_{\mathrm U}
\hat r$ is exactly \eqref{eq:TV}. ULS needs $\Gamma$ only for the reference
distribution, never to compute the estimator or the statistic; it is the
cheapest and most robust choice, and it inherits nothing from the model's
conditioning.

\paragraph{NTML, as reweighted least squares.} With
\[
  F_{\mathrm M}(\theta)=\log|\Sigma(\theta)|
    +\tr\{\widehat\Sigma\,\Sigma(\theta)\Ninv\}-\log|\widehat\Sigma|-p,
\]
a second-order expansion about exact fit gives
\begin{equation}\label{eq:rls}
  T_{\mathrm M}=2NF_{\mathrm M}(\hat\theta)
    =N\hat r\T V_{\mathrm M}\hat r+o_p(1),\qquad
  V_{\mathrm M}=\tfrac12 D\T(\Sigma\Ninv\otimes\Sigma\Ninv)D,
\end{equation}
the \emph{reweighted least squares} (RLS) quadratic form of Browne
\cite{browne1974}: the ML statistic linearises to a residual quadratic form with
the model-implied normal-theory weight. The weight $V_{\mathrm M}$ is evaluated
at $\Sigma(\hat\theta)$, hence a smooth function of the fitted model, not of the
data. This is exactly why NTML is admissible here and GLS is not: GLS would use
$\widehat S\Ninv\otimes\widehat S\Ninv$, a sample weight.

\begin{remark}
For the \emph{statistic} a data weight is in fact asymptotically negligible
(it multiplies two $O_p(N^{-1/2})$ residual factors), so GLS would give a valid
test too. We drop it only to keep the weight a function of the model and to keep
the normal-theory reductions of \S\ref{sec:normal} clean; it adds nothing this
project needs.
\end{remark}

Both $T_{\mathrm U}$ and $T_{\mathrm M}$ are instances of \eqref{eq:TV} with
$M=I-\Delta J_G$ (the estimator's projector), so Proposition~\ref{prop:gof}
governs both.

\section{Reductions under normality}\label{sec:normal}

Substitute $\Gamma=\Gamma_{\mathrm N}$. The normal-theory weight is a
generalised inverse of the normal-theory moment covariance on the symmetric
subspace,
\begin{equation}\label{eq:VM-ginv}
  V_{\mathrm M}\,\Gamma_{\mathrm N}\,V_{\mathrm M}=V_{\mathrm M},
\end{equation}
so $V_{\mathrm M}$ plays the role of $\Gamma_{\mathrm N}^{+}$. This is what makes
the NTML statistic reduce, and it isolates the efficiency loss cleanly.

\begin{proposition}[Normal-theory NTML inflation]\label{prop:inflation}
Under normality, with the NTML statistic and any Fisher-consistent $\tau$,
\begin{equation}\label{eq:c-ge-1}
  \E\,T_{\mathrm M}\to
  \tr\!\big(V_{\mathrm M}M\Gamma_{\mathrm N}M\T\big)\ge df,
\end{equation}
with equality iff $M$ is the NTML (efficient) projector $M_{V_{\mathrm M}}$, i.e.
iff $J_G$ coincides with the ML Jacobian. Hence the Satorra--Bentler constant
$c=\E\,T_{\mathrm M}/df\ge 1$: the raw NTML statistic evaluated at the
non-iterative estimates is inflated relative to $\chi^2_{df}$ and must be scaled
\emph{down}.
\end{proposition}

\begin{proof}[Sketch]
Using \eqref{eq:VM-ginv}, $\tr(V_{\mathrm M}M\Gamma_{\mathrm N}M\T)$ is the
generalised residual variance of a projector $M$ onto the complement of
$\Delta$; among all such projectors it is minimised, at value $df$, by the
$V_{\mathrm M}$-orthogonal projector $M_{V_{\mathrm M}}$ (the Gauss--Markov /
optimal-GLS residual argument). Any other consistent $\tau$ gives $M\neq
M_{V_{\mathrm M}}$ and a strictly larger trace.
\end{proof}

This is the precise content of ``the statistic reduces under normality but the
distribution does not.'' Under normality $T_{\mathrm M}$ collapses to the RLS
form \eqref{eq:rls} and its scaling constant is the explicit
\eqref{eq:c-ge-1}; but the reference law stays a weighted sum
\eqref{eq:weighted-chi2} with $\lambda_j\not\equiv1$, and only an efficient
(hence iterative) estimator would return $\chi^2_{df}$. For ULS the departure is
even plainer: $V_{\mathrm U}\neq\Gamma_{\mathrm N}^{+}$, so even the ULS
\emph{minimiser} is a weighted sum of $\chi^2$ under normality (the known ULS
fact); the non-iterative $\tau$ only changes which weights.

\section{Nested tests}\label{sec:nested}

Let $M_0\subset M_1$ be nested by a smooth restriction $a(\theta)=0$,
$a:\mathbb R^{q_1}\to\mathbb R^{df_d}$, $df_d=q_1-q_0$, with
$A=\partial a/\partial\theta\T$ of full row rank at $\theta_0$.

\paragraph{Wald (recommended).} Fit the larger model with its non-iterative map,
$\hat\theta_1=\tau_1(\hat\sigma)$, covariance
$\widehat\Omega_1=J_1\widehat\Gamma J_1\T/N$ (the delta-method object from the SE
note). Then
\begin{equation}\label{eq:wald}
  W=N\,a(\hat\theta_1)\T\big(A\widehat\Omega_1 A\T\big)\Ninv a(\hat\theta_1)
  \Rightarrow\chi^2_{df_d}
\end{equation}
\emph{exactly}, with no eigenvalue spread and no second $\Gamma$-quadratic. It
reuses only machinery already built for standard errors and is the natural
nested test for these estimators; efficiency of $\tau$ costs power, not
validity.

\paragraph{Difference (pseudo-LRT).} If an LRT-shaped statistic is wanted,
\begin{equation}\label{eq:Td}
  T_d=T_{V,0}-T_{V,1}
    =2N\{F_V(\hat\theta_0)-F_V(\hat\theta_1)\}
  \Rightarrow\sum_j\omega_j\chi^2_{1,j},
\end{equation}
with $\omega_j$ the nonzero eigenvalues of $U_d\Gamma$,
$U_d=M_0\T V M_0-M_1\T V M_1$, scaled Satorra-2000 style by
$c_d=\tr(U_d\Gamma)/df_d$ \cite{satorra2000,satorrabentler2001}. Two caveats
distinguish this from the ML case:
\begin{enumerate}
\item Because $\hat\theta_0,\hat\theta_1$ do not minimise $F_V$, $U_d$ is
  generally indefinite, so $T_d$ can be \emph{negative} in finite samples (the
  monotonicity $T_{V,0}\ge T_{V,1}$ that holds for minimisers is not guaranteed).
\item $\rank U_d$ need not equal $df_d$; the mean $\tr(U_d\Gamma)$ is still
  positive and $\approx df_d\,c$, so mean scaling remains well defined.
\end{enumerate}
The Wald form \eqref{eq:wald} avoids both. We keep \eqref{eq:Td} for the
ML-analogy and because, specialised to a discrepancy-minimising $\tau$, it is
exactly the standard scaled difference test.

\section{Robust default, categorical outlook, implementation}\label{sec:impl}

Everything above is written for a general $\Gamma$. Taking $\widehat\Gamma$
empirical yields distribution-free GOF, Wald, and difference tests; taking
$\Gamma_{\mathrm N}$ gives the normal-theory specialisations of
\S\ref{sec:normal}. The estimator itself never touches $\Gamma$; only the tests
do. The motivating extension is categorical data: a non-iterative estimator on
polychoric moments with $\Gamma$ the polychoric ADF covariance drops into the
same three formulas unchanged, which is where a convergence-free estimator plus
a robust test is most useful and least served by current tooling.

The pieces already exist in \texttt{magmaan}: $\Delta$ from
\texttt{ModelEvaluator}; $J_G$ from \texttt{src/estimate/start\_guttman.cpp}
and the companion note (finite-difference validated to $\sim1\%$ in
experiment~40); $\Gamma$ in both normal-theory and empirical forms; the
Satorra--Bentler / pEBA spectrum reducer in \texttt{robust}; and the Wald
machinery. The only new object is $M=I-\Delta J_G$ and the two weights
$V_{\mathrm U}=D\T D$, $V_{\mathrm M}=\tfrac12 D\T(\Sigma\Ninv\otimes
\Sigma\Ninv)D$; feed $M\T V M$ times $\Gamma$ to the existing reducer.
Implementation is deferred.

\begin{thebibliography}{9}

\bibitem{bollen1996}
Bollen, K. A. (1996).
An alternative two stage least squares (2SLS) estimator for latent variable
equations.
\emph{Psychometrika}, 61, 109--121.
\href{https://doi.org/10.1007/BF02296961}{doi:10.1007/BF02296961}.

\bibitem{browne1974}
Browne, M. W. (1974).
Generalized least squares estimators in the analysis of covariance structures.
\emph{South African Statistical Journal}, 8, 1--24.

\bibitem{browne1984}
Browne, M. W. (1984).
Asymptotically distribution-free methods for the analysis of covariance
structures.
\emph{British Journal of Mathematical and Statistical Psychology}, 37, 62--83.
\href{https://doi.org/10.1111/j.2044-8317.1984.tb00789.x}
{doi:10.1111/j.2044-8317.1984.tb00789.x}.

\bibitem{dhaene2024}
Dhaene, S., \& Rosseel, Y. (2024).
An evaluation of non-iterative estimators in confirmatory factor analysis.
\emph{Structural Equation Modeling}, 31(1), 1--13.
\href{https://doi.org/10.1080/10705511.2023.2187285}
{doi:10.1080/10705511.2023.2187285}.

\bibitem{guttman1952}
Guttman, L. (1952).
Multiple group methods for common-factor analysis: Their basis, computation, and
interpretation.
\emph{Psychometrika}, 17, 209--222.
\href{https://doi.org/10.1007/BF02288783}{doi:10.1007/BF02288783}.

\bibitem{hagglund1982}
H\"agglund, G. (1982).
Factor analysis by instrumental variables methods.
\emph{Psychometrika}, 47, 209--222.
\href{https://doi.org/10.1007/BF02296273}{doi:10.1007/BF02296273}.

\bibitem{satorra1989}
Satorra, A. (1989).
Alternative test criteria in covariance structure analysis: A unified approach.
\emph{Psychometrika}, 54, 131--151.
\href{https://doi.org/10.1007/BF02294453}{doi:10.1007/BF02294453}.

\bibitem{satorra2000}
Satorra, A. (2000).
Scaled and adjusted restricted tests in multi-sample analysis of moment
structures.
In R. D. H. Heijmans, D. S. G. Pollock, \& A. Satorra (Eds.),
\emph{Innovations in multivariate statistical analysis} (pp. 233--247). Kluwer.

\bibitem{satorrabentler1994}
Satorra, A., \& Bentler, P. M. (1994).
Corrections to test statistics and standard errors in covariance structure
analysis.
In A. von Eye \& C. C. Clogg (Eds.),
\emph{Latent variables analysis} (pp. 399--419). Sage.

\bibitem{satorrabentler2001}
Satorra, A., \& Bentler, P. M. (2001).
A scaled difference chi-square test statistic for moment structure analysis.
\emph{Psychometrika}, 66, 507--514.
\href{https://doi.org/10.1007/BF02296192}{doi:10.1007/BF02296192}.

\bibitem{shapiro1986}
Shapiro, A. (1986).
Asymptotic theory of overparameterized structural models.
\emph{Journal of the American Statistical Association}, 81, 142--149.
\href{https://doi.org/10.1080/01621459.1986.10478250}
{doi:10.1080/01621459.1986.10478250}.

\end{thebibliography}

\end{document}
